\section{Deployment and a relative cost lower bound}
\label{deployment:section}

We next turn a graphing of an amalgamated product action into a graphing
of one factor with the amalgamating relation supplied for free.
The construction is Gaboriau's deployment of a graphing
\cite[\S\,IV.28--IV.32]{gaboriau-2000}, specialized to finitely
many sheets. The compression lemma accounts for its cost.
Throughout this section,
\(\Gamma=A*_{J}B\) is finitely generated, \(J\) is infinite, and
\(\Gamma\) acts essentially freely by measure-preserving Borel
automorphisms on a standard probability space \((X,\mu)\).
We use right actions and write \(\mathcal R_L\) for the
orbit relation of a subgroup \(L\leq\Gamma\).  Define
\[
 \delta_X=\inf\bigl\{C(\mathcal E):
   \mathcal E\text{ is a graphing in }\mathcal R_A,
   \quad\mathcal R_J\vee\mathcal E=\mathcal R_A\bigr\}.
\]
Thus only the edges added to \(\mathcal R_J\) are charged in
\(\delta_X\).

\begin{proposition}
\label{deployment:main}
Let $\Gamma=A*_J B$ be finitely generated, let $J$ be infinite, and
let $\Gamma$ act essentially freely and p.m.p. on a standard probability
space $(X,\mu)$. If $\delta_X$ is the infimum cost of graphings that
generate $\mathcal R_A$ together with the supplied relation
$\mathcal R_J$, then
\[
 \operatorname{Cost}(\mathcal R_\Gamma)\geq 1+\delta_X.
\]
\end{proposition}

We prove the proposition using graphings at finite stages.  The only
limit will concern the sizes of increasing equivalence classes.
We first pass to an invariant conull Borel set on which the action is
free.  Generation identities may also be made exact after removing
countably many null sets and their saturations under the ambient
countable p.m.p. relation.  These restrictions leave all costs unchanged.

\subsection{A finite graphing and its subdivision}
\label{deployment:finite}

Fix \(\rho>0\).  There is a finite generating graphing of
\(\mathcal R_\Gamma\), consisting of restrictions of nonidentity group
translations, whose cost \(c\) satisfies
\begin{equation}
\label{deployment:near-optimal}
 c<\operatorname{Cost}(\mathcal R_\Gamma)+\rho.
\end{equation}
Here is the approximation argument.  Start with a countable generating
graphing of cost less than
\(\operatorname{Cost}(\mathcal R_\Gamma)+\rho/2\), and split each map
into disjoint Borel pieces according to its group displacement.  This
does not change its cost; discard identity restrictions, which can only
lower it.  Let \(\Psi_N\) be the first \(N\) maps of
the resulting graphing and let \(T_N\) be their generated relation.
For a fixed finite generating set \(K\) of \(\Gamma\) not containing
the identity, put
\[
 D_{N,s}=\{x\in X:(x,xs)\notin T_N\},\qquad s\in K.
\]
Each set is Borel because membership in \(T_N\) is witnessed by a
finite path.  For each \(s\), these sets decrease to a null set, so
\(\sum_{s\in K}\mu(D_{N,s})\to0\).  Adding the translation by \(s\)
on \(D_{N,s}\), for every \(s\in K\), gives a finite generating
graphing and proves \eqref{deployment:near-optimal} for large \(N\).

Write its maps as \(x\mapsto xg_h\) on Borel domains \(D_h\), with
\(h\) in a finite index set.  Choose a factor
word
\[
 g_h=g_{h,1}\cdots g_{h,\ell_h},\qquad
 g_{h,i}\in A\cup B,\quad \ell_h\geq1.
\]
For letters in \(J\), choose either factor type.  Form a space \(Z\)
by adjoining \(\ell_h-1\) disjoint copies of \(D_h\) to \(X\), one
for each internal vertex of this word.  Give each copy its inherited
measure and denote the resulting, unnormalized measure by \(\lambda\).
Then
\[
 \lambda(Z)=1+\sum_h(\ell_h-1)\mu(D_h)<\infty.
\]
Define \(P:Z\to X\) to be the identity on \(X\), and on the copy
indexed by \((h,i)\) define
\[
 P(x,h,i)=xg_{h,1}\cdots g_{h,i}.
\]
The restriction of \(P\) to each sheet is a measure-preserving Borel
bijection onto its image.  For \(L=\Gamma,A,B,J\), let
\[
 \overline{\mathcal R}_L
   =\{(z,z')\in Z^2:(Pz,Pz')\in\mathcal R_L\}.
\]
These are countable p.m.p. Borel equivalence relations.  Indeed, if
\(P_r\) and \(P_s\) denote the restrictions to two sheets, the partial
maps \(P_s^{-1}\circ(x\mapsto x\ell)\circ P_r\), for \(\ell\in L\),
enumerate their pairs and preserve measure on their natural domains.

Replace each original edge by the chain through its new vertices, and
divide the successive edges according to their factor types into
\(\mathcal E_A^0\) and \(\mathcal E_B^0\).  Every chain edge belonging
to the \(h\)-th map has cost \(\mu(D_h)\), so
\begin{equation}
\label{deployment:initial-cost}
 C(\mathcal E_A^0)+C(\mathcal E_B^0)
   =c+\lambda(Z)-1.
\end{equation}
Their join is \(\overline{\mathcal R}_\Gamma\).  To verify this,
first connect base-sheet points by replacing paths in the original
graphing by their subdivided paths.  Every additional point is attached
along its chain to a base-sheet point in the same projected orbit.
This also connects arbitrary pairs in the pullback relation.

\subsection{Increasing the supplied relation}
\label{deployment:iteration}

Set \(S_0=\mathrm{id}_Z\).  We construct relations
\(S_n\subseteq\overline{\mathcal R}_J\) and graphings
\(\mathcal E_i^n\subseteq\overline{\mathcal R}_i\), for \(i=A,B\),
and put
\[
 Q_i^n=S_n\vee\mathcal E_i^n.
\]
Fix \(\epsilon>0\), and choose positive errors \(e_n\) with
\(\sum_{n\geq0}e_n<\epsilon\).
At stage \(n\), choose \(i=A\) or \(B\) alternately and set
\[
 S_{n+1}=Q_i^n\cap\overline{\mathcal R}_J.
\]
We have \(S_n\subseteq S_{n+1}\subseteq Q_i^n\).  Apply
Lemma~\ref{compression:main} with
\(S=S_n\), \(T=S_{n+1}\), \(Q=Q_i^n\), and error \(e_n\).
It replaces \(\mathcal E_i^n\) by \(\mathcal E_i^{n+1}\) such that
\[
 Q_i^n=S_{n+1}\vee\mathcal E_i^{n+1}
\]
and
\[
 C(\mathcal E_i^{n+1})
 \leq C(\mathcal E_i^n)-\kappa(S_n)+\kappa(S_{n+1})+e_n.
\]
Leave the other edge family unchanged.  The processed relation
\(Q_i^{n+1}\) equals \(Q_i^n\); the other relation grows because its
supplied relation grows.  Thus both sequences \(Q_i^n\) increase,
remain inside their respective factor relations, and have join
\(\overline{\mathcal R}_\Gamma\) at every stage.  The edge families
themselves need not be increasing.

All costs and all values of \(\kappa\) here use \(\lambda\), without
normalizing its total mass.  Summing the compression inequalities and
using \(\kappa(S_0)=\lambda(Z)\) and
\eqref{deployment:initial-cost} gives
\begin{equation}
\label{deployment:finite-bound}
 C(\mathcal E_A^n)+C(\mathcal E_B^n)
  \leq c-1+\kappa(S_n)+\epsilon
  \qquad(n\geq0).
\end{equation}
In particular all graphings used in subsequent applications of the
compression lemma have finite cost.

\subsection{Exhausting the amalgamating relation}
\label{deployment:limit}

Let \(S_\infty=\bigcup_nS_n\) and
\(Q_i^\infty=\bigcup_nQ_i^n\).  These are Borel equivalence relations,
and their monotonicity gives
\begin{equation}
\label{deployment:intersections}
 Q_A^\infty\vee Q_B^\infty=\overline{\mathcal R}_\Gamma,
 \qquad Q_i^\infty\cap\overline{\mathcal R}_J=S_\infty.
\end{equation}
For the nontrivial inclusion in the second identity, a pair in
\(Q_i^n\cap\overline{\mathcal R}_J\) enters \(S_{m+1}\) at the next
later stage \(m\) at which type \(i\) is processed.

We claim that \(S_\infty=\overline{\mathcal R}_J\).  Take a pair in
\(\overline{\mathcal R}_J\) and a shortest chain between its endpoints
with steps in \(Q_A^\infty\) or \(Q_B^\infty\).  If the chain has at
least two steps, its types alternate: adjacent steps of the same type
could be composed.  Moreover, no step lies in
\(\overline{\mathcal R}_J\).  By
\eqref{deployment:intersections}, such a step would belong to
\(S_\infty\), hence to both factor relations, and could be composed
with a neighboring step.

Projecting the chain to \(X\), freeness gives unique group
displacements for its steps.  These form an alternating word with
letters outside \(J\), whose product belongs to \(J\) because the
projected endpoints are \(J\)-related.  This contradicts amalgam
normal form.  The shortest chain therefore has length at most one;
\eqref{deployment:intersections} puts its endpoint pair in
\(S_\infty\), as claimed.

Every \(\overline{\mathcal R}_J\)-class is infinite: the class of
\(z\) contains the distinct base-sheet points \(P(z)j\), for
\(j\in J\).  Since \([z]_{S_n}\) increases with this infinite union,
\[
 \frac{1}{|[z]_{S_n}|}\downarrow0.
\]
The integrands are bounded by \(1\) on the fixed finite-measure space
\(Z\).  Dominated convergence consequently gives
\begin{equation}
\label{deployment:kappa-limit}
 \kappa(S_n)\longrightarrow0.
\end{equation}

\subsection{Returning to the original space}
\label{deployment:projection}

Fix a finite stage \(n\).  For each map \(\varphi\) in
\(\mathcal E_i^n\), partition its domain according to both its source
sheet \(r\) and its target sheet \(s\).  On such a piece \(D\), the
map
\[
 P_s\circ\varphi\circ P_r^{-1}:
 P_r(D)\longrightarrow P_s(\varphi(D))
\]
is a p.m.p. Borel partial isomorphism on \(X\), belongs to
\(\mathcal R_i\), and has cost \(\lambda(D)\).  These maps form a
graphing \(\mathcal F_i^n\) with
\begin{equation}
\label{deployment:projection-cost}
 C(\mathcal F_i^n)\leq C(\mathcal E_i^n).
\end{equation}
This construction uses the sheet restrictions of \(P\); it requires
no injectivity of \(P\) on all of \(Z\).

Put \(T_i=\mathcal R_J\vee\mathcal F_i^n\).  Projecting paths from
the base sheet shows that
\(T_A\vee T_B=\mathcal R_\Gamma\): every \(S_n\)-step projects into
\(\mathcal R_J\), and all remaining steps project into the indicated
graphings.  We show that \(T_A=\mathcal R_A\).
For a pair in \(\mathcal R_A\), choose a shortest chain with steps
from \(T_A,T_B\).  A chain of length at least two would alternate and
contain no \(\mathcal R_J\)-step, since that relation is contained in
both \(T_A\) and \(T_B\).  Its displacement word would alternate
outside \(J\), have length at least two, and belong to \(A\), again
contradicting amalgam normal form.  A single step of type \(B\)
between \(A\)-related endpoints has displacement in
\(A\cap B=J\), so it already belongs to \(T_A\).  This proves the
claim.

Thus \(\mathcal F_A^n\) is admissible in the definition of
\(\delta_X\).  Equations \eqref{deployment:finite-bound} and
\eqref{deployment:projection-cost} give, at every finite stage,
\[
 \delta_X\leq C(\mathcal F_A^n)
 \leq C(\mathcal E_A^n)+C(\mathcal E_B^n)
 \leq c-1+\kappa(S_n)+\epsilon.
\]
Let \(n\to\infty\) using \eqref{deployment:kappa-limit}, then let
\(\epsilon\) and \(\rho\) tend to zero in
\eqref{deployment:near-optimal}.  This proves
Proposition~\ref{deployment:main}.
